\section{M5 --- Avaliação Quantitativa por Injeção Sintética}
\label{sec:m5}

\subsection{Motivação e desenho}
O problema é não supervisionado: não há rótulos de anomalia reais. Para obter
métricas quantitativas (\emph{Precision}, \emph{Recall}, F1) injetam-se anomalias em
posições \emph{conhecidas} da série de teste, produzindo um \emph{ground-truth}
controlado (ADR-0006). Foram inseridos $n=50$ \emph{price shocks} por ativo, em posições
aleatórias (\emph{seed} fixa), somando um delta de magnitude \texttt{price\_shock}$=0{,}15$
ao log-retorno \textbf{já escalado} --- o mesmo espaço em que o modelo foi treinado.
Injetar no preço bruto produziria choques de magnitude variável conforme o nível de preço
do ativo.

\subsection{Regra de rotulagem por janela}
Como o modelo processa janelas de 30 passos, um choque em qualquer posição eleva o erro de
reconstrução de toda a janela. Adota-se, portanto, a regra: uma janela é positiva se
\textbf{qualquer} passo injetado cair dentro dela. Os vetores de rótulo e de \emph{flags}
do detector ficam alinhados janela a janela, seguindo a mesma indexação de
\texttt{make\_windows} (Seção~\ref{sec:m2}), o que evita desalinhamento temporal no cálculo
das métricas.

\subsection{Resultados}
A Tabela~\ref{tab:metrics} apresenta as métricas por ativo, comparando limiar estático
(p95 do treino) e dinâmico (calculado sobre os erros da série perturbada). O padrão é
claro e metodologicamente informativo:
\begin{itemize}[noitemsep]
  \item O \emph{recall} é \textbf{uniformemente baixo} nos ativos de regime estável
        (PETR4, VALE3: recall $<7\%$). Com magnitude \emph{absoluta} de $0{,}15$ no espaço
        escalado, o choque sintético é comparável ao ruído normal desses ativos e raramente
        empurra a janela acima do limiar.
  \item Onde o regime de erro de base é mais alto (AMER3, ITUB4), os choques são detectados
        com mais frequência (AMER3 estático: recall $\approx$30\%, F1 $\approx$0{,}42).
  \item O limiar dinâmico tende a elevar a \emph{precision} (ex.\ VALE3 e ITUB4 acima de
        0{,}9), por se ajustar ao ruído local, mas sem ganho de \emph{recall}.
\end{itemize}

\begin{table}[h]
\centering
\caption{Precision / Recall / F1 do detector contra o \emph{ground-truth} sintético
($n=50$ \emph{price shocks}, \texttt{price\_shock}$=0{,}15$).}
\label{tab:metrics}
\begin{tabular}{llccc}
\toprule
\textbf{Ticker} & \textbf{Limiar} & \textbf{Precision} & \textbf{Recall} & \textbf{F1} \\
\midrule
PETR4.SA & estático & 0{,}139 & 0{,}006 & 0{,}011 \\
PETR4.SA & dinâmico & 0{,}548 & 0{,}020 & 0{,}038 \\
VALE3.SA & estático & 0{,}111 & 0{,}005 & 0{,}009 \\
VALE3.SA & dinâmico & 0{,}964 & 0{,}061 & 0{,}115 \\
AMER3.SA & estático & 0{,}694 & 0{,}303 & 0{,}421 \\
AMER3.SA & dinâmico & 0{,}605 & 0{,}120 & 0{,}200 \\
ITUB4.SA & estático & 0{,}769 & 0{,}130 & 0{,}222 \\
ITUB4.SA & dinâmico & 0{,}939 & 0{,}036 & 0{,}069 \\
\bottomrule
\end{tabular}
\end{table}

\subsection{Leitura crítica}
O \emph{recall} baixo \textbf{não} indica falha do detector, mas que a magnitude de injeção é
\textbf{provisória e absoluta}: a dificuldade da tarefa varia arbitrariamente entre ativos,
penalizando os de baixa volatilidade. Esta é a evidência empírica que motiva a recomendação
do ADR-0006: calibrar o choque em $k$ desvios-padrão do retorno \emph{de cada ativo}, tornando
a dificuldade comparável entre setores. A avaliação corrente, portanto, é um piso
metodológico --- valida o \emph{pipeline} de métricas (injeção $\to$ rotulagem $\to$
P/R/F1), e não um resultado final. Adicionalmente, a taxa-base de $\approx$5\% imposta pelo
p95 (Seção~\ref{sec:m4}) limita a \emph{precision} máxima atingível e deve ser lida em
conjunto. A injeção de \emph{volume spikes} (chave \texttt{volume\_spike} no
\texttt{config.yaml}) permanece prevista mas ainda não implementada nesta etapa.
